\documentclass[12pt]{article}
\usepackage[slovene]{babel}
\usepackage{printlen}

\usepackage{ragged2e,lmodern} %sodobni LaTeX

\usepackage{array,booktabs, multirow} % Tabele
\usepackage{graphicx} % Slike
\usepackage{tcolorbox}
\tcbset{colback=white, arc=0mm, boxrule=0.2mm,}

\usepackage{mathtools}
\usepackage{amsthm, amssymb} % Matematik
\usepackage{lipsum}

% \usepackage{hyperref} % Hiperpovezave

\usepackage{minted}


\title{\LaTeX{} 4. Matematika}
\author{Antonio Montero}
\date{22. oktober 2025}



\newsavebox{\mybox}



\newcommand{\pbox}[1]{
  \savebox{\mybox}{\tcbox{#1}}
  
  \usebox{\mybox}
  
  \vspace{1cm}
  \the\wd\mybox
  \hrule 
  \vspace{1cm}
  }
  



  \NewDocumentCommand{\kpi}{}{\mathrm{\pi}}


  \theoremstyle{plain}
  \newtheorem{izrek}{Izrek}
  \newtheorem{trd}[izrek]{Trditev}
  \theoremstyle{definition}
  \newtheorem{defja}[izrek]{Definicija}
  \theoremstyle{remark}
  \newtheorem*{zgled}{Zgled}

\begin{document}
\maketitle

\begin{tcolorbox}
  % \begin{equation}
  % \label{eq:kvadratna_formula}
  % % \tag{Kvadratna formula}
  %   x_{1,2} = \frac{-b \pm \sqrt{b^2 - 4ac}}{2a}
  % \end{equation}

V enačbi \ref{eq:kvadratna_formula} je prikazana kvadratna formula.
\end{tcolorbox}

\begin{tcolorbox}
\[2+ \text{enka} = 3 \]
\end{tcolorbox}

\vspace{2cm}

\pbox{$\displaystyle 2+ \text{enka} = 3 $}
\pbox{%
$123xyz$ je enako kot \(1 2 3 x y z \)
}

\pbox{
  Primerjaj besedo \textit{office} z besedo $office$.
}

$\kpi$

$\begin{smallmatrix}
a & b \\ c & d
\end{smallmatrix}$

\[\sum_{j\in\{0,\ldots\, 10\}}
\sum_{i\in\{0,\ldots\, 3^{3^{3^j}}\}} i\cdot j\]

$\mathrm{e}$


\vspace{2cm}

% \begin{izrek}[Pitagorov izrek]
%   Vsota površin kvadratov katet pravokotnega 
%   trikotnika je enaka površini kvadrata nad 
%   hipotenuzo.
%   Izrek lahko zapišemo tudi kot: \[c^{2} = a^{2} + b^{2}\]
% \end{izrek}

\vspace{2cm}

V \LaTeX-u moramo matematični del besedila zapisati v matematičnem načinu. Tudi če je to nekaj tako preprostega kot ta številka $2$ ali ta spremenljivka $x$.

\vspace{2cm}

\section{Razdelek z teoremi}

% \begin{izrek} (Pitagorov izrek)
%   V pravokotnem trikotniku velja \[c^{2} = a^{2} + b^{2}\] 
% \end{izrek}
% % ...

\vspace{2cm}
\begin{trd} 
  Ne vem, če ta trditev velja.
\end{trd}
\begin{proof}
  Očitno, uporabimo
  \[E=mc^2.\qedhere\]
\end{proof}
% % ... 
% \begin{defja}
%   Tukaj sem napisal definicijo. Običajno vporabljamo \emph{povdarjeno} besedilo za naslove definicij.
% \end{defja}
% % ... 
% \begin{zgled}
%   Tukaj sem napisal primer.
% \end{zgled}


\clearpage

% Maxwellove enačbe v vakuumu so:
% \begin{align}
%   \nabla \cdot \mathbf{E} & = \frac{\rho}{\varepsilon_0} \label{eq:GaussE} \\
%   \nabla \cdot \mathbf{B} & = 0 \label{eq:GaussB} \\
%   \nabla \times \mathbf{E} & = -\frac{\partial \mathbf{B}}{\partial t} \label{eq:Faraday} \\
%   \nabla \times \mathbf{B} & = \mu_0 \mathbf{J} + \mu_0 \varepsilon_0 \frac{\partial \mathbf{E}}{\partial t} \label{eq:Ampere}
% \end{align} 

Enačba \ref{eq:GaussE} opisuje Gaussov zakon za električno polje, enačba \ref{eq:GaussB} pa Gaussov zakon za magnetno polje. Enačba \ref{eq:Faraday} predstavlja Faradayev zakon elektromagnetne indukcije, medtem ko enačba \ref{eq:Ampere} predstavlja Amperov zakon z Maxwellovim dodatkom.

\clearpage

\begin{izrek}[de Morganovi zakoni]
  Za vsakih dvoje podmnožic $A$ in $B$ množice $X$ veljata naslednji enačbi:
  \begin{align}
    X\setminus (A \cup B) & = (X \setminus A) \cap (X \setminus B) \label{eq:deMorgan1}\\
    X\setminus (A \cap B) & = (X \setminus A) \cup (X \setminus B) \label{eq:deMorgan2}
  \end{align}
\end{izrek}
\begin{proof}
  Dokaz za prvo enačbo:
  Naj bo $x \in X\setminus (A \cup B)$. Potem 
  \begin{align*}
    x \in X \text{ in } x \notin A \cup B
    & \Rightarrow x \in X \text{ in } (x \notin A \text{ in } x \notin B) \\
    & \Rightarrow  (x \in X \text{ in } x \notin A) \text{ in } (x \in X \text{ in } x \notin B) \\
    & \Rightarrow x \in X \setminus A \text{ in } x \in X \setminus B \\
    & \Rightarrow  x \in (X \setminus A) \cap (X \setminus B)
  \end{align*}
%
  Obratno, naj bo $x \in (X \setminus A) \cap (X \setminus B)$. Potem 
  \begin{align*}
    x \in (X \setminus A) \cap (X \setminus B)
    & \Rightarrow (x \in X \text{ in } x \notin A) \text{ in } (x \in X \text{ in } x \notin B) \\
    & \Rightarrow x \in X \text{ in } (x \notin A \text{ in } x \notin B) \\
    & \Rightarrow x \in X \text{ in } x \notin A \cup B \\
    & \Rightarrow x \in X\setminus (A \cup B)
  \end{align*}
S tem je dokaz prve enačbe zaključen.

  Dokaz za drugo enačbo je podoben: 
  Naj bo $x \in X\setminus (A \cap B)$. Potem 
  \begin{align*}
    x \in X\setminus (A \cap B)
    & \Rightarrow x \in X \text{ in } x \notin A \cap B \\
    & \Rightarrow x \in X \text{ in } (x \notin A \text{ ali } x \notin B) \\
    & \Rightarrow  (x \in X \text{ in } x \notin A) \text{ ali } (x \in X \text{ in } x \notin B) \\
    & \Rightarrow x \in X \setminus A \text{ ali } x \in X \setminus B \\
    & \Rightarrow  x \in (X \setminus A) \cup (X \setminus B)
  \end{align*}
%  
Obratno, naj bo $x \in (X \setminus A) \cup X \setminus B$. Potem 
  \begin{align*}
    x \in (X \setminus A) \cup (X \setminus B)
    & \Rightarrow (x \in X \text{ in } x \notin A) \text{ ali } (x \in X \text{ in } x \notin B) \\
    & \Rightarrow x \in X \text{ in } (x \notin A \text{ ali } x \notin B) \\
    & \Rightarrow x \in X \text{ in } x \notin A \cap B \\
    & \Rightarrow x \in X\setminus (A \cap B)
  \end{align*}
  Sledi, da $X\setminus (A\cap B) = (X \setminus A) \cup (X \setminus B)$.

  Torej enačbi \ref{eq:deMorgan1} in \ref{eq:deMorgan2} držita.
\end{proof}

\vspace{2cm}

\begin{minipage}{0.5\textwidth}
\begin{multline*}
  a + b + c + d + e \\
  + f + g + h + i \\
  = j + k + l + m + n
  \end{multline*}
\end{minipage}

\vspace{2cm}

\begin{minipage}{0.5\textwidth}
\begin{multline*}
    a + b + c \\
    \shoveleft{+ d + e + f} \\
    \shoveright{+ g + h + i} \\
    = j + k + l + m + n
\end{multline*}
\end{minipage}

\vspace{2cm}

\begin{minipage}{0.5\textwidth}
  \begin{multline*}
    f(a, b) = \left(
    \int_a^b f(x)
    dooooooooolgo
    \right. \\
    \left.
    \vphantom{
    \int_a^b f(x)
    dooooooooolgo
    }
    kratko d x
    \right)
  \end{multline*}
\end{minipage}



\end{document}